\section{Paradigma Estatístico: Cox PWP-GT}

O modelo de Prentice-Williams-Peterson será a nossa âncora estatística para a recorrência.
\begin{itemize}
    \item \textbf{Estratificação:} O modelo será estratificado pelo número de episódios (\textit{strata}), permitindo que o hazard base varie em cada recidiva.
    \item \textbf{Variáveis Dependentes do Tempo:} Incorporação de indicadores macro dinâmicos (PIB, Inflação, Juros) como covariáveis que evoluem em cada intervalo.
\end{itemize}

\section{Ensemble Learning: Random Survival Forests (RSF)}

O RSF é utilizado para capturar as complexas interações não-lineares entre rácios financeiros e indicadores de \textit{Governance}.
\begin{itemize}
    \item \textbf{Splitting Rule:} Utilização do \textit{Log-rank} para maximizar a diferença estatística entre as curvas de sobrevivência dos nós filhos.
    \item \textbf{Interpretabilidade:} Aplicação de valores SHAP (\textit{SHapley Additive exPlanations}) para decompor a contribuição de cada rácio para o risco individual da PME.
\end{itemize}

\section{Deep Learning: DynamicDeepHit}

Este modelo representa a inovação central do projeto, ao introduzir a memória sequencial na análise de sobrevivência.
\begin{itemize}
    \item \textbf{Encoder LSTM:} Processa a sequência temporal de rácios financeiros dos últimos 5 anos, extraindo um \textit{embedding} que representa a trajetória de saúde da empresa.
    \item \textbf{Decoders Causa-Específicos:} Redes neuronais independentes para prever a probabilidade de transição para Distress (Recorrência) vs. Falência (Risco Concorrente).
    \item \textbf{Função de Perda Híbrida:} Combinação da \textit{Log-Likelihood} para o tempo do evento com uma \textit{Ranking Loss} para garantir a ordenação correta do risco.
\end{itemize}
